\documentclass[10pt,a4paper]{amsart}
\usepackage{amsmath,amssymb,amsthm,graphicx}
% For commutative diagrams we use Xy-pic
\usepackage[all]{xy}

% We use 2cell for 2-commutative diagrams.
\xyoption{2cell}
% Theorem environments.
%
\theoremstyle{plain}
\newtheorem{theorem}[subsection]{Theorem}
\newtheorem{proposition}[subsection]{Proposition}
\newtheorem{lemma}[subsection]{Lemma}

\theoremstyle{definition}
\newtheorem{definition}[subsection]{Definition}
\newtheorem{example}[subsection]{Example}
\newtheorem{exercise}[subsection]{Exercise}
\newtheorem{situation}[subsection]{Situation}

\theoremstyle{remark}
\newtheorem{remark}[subsection]{Remark}
\newtheorem{remarks}[subsection]{Remarks}

\numberwithin{equation}{subsection}

% Macros
%
\def\lim{\mathop{\mathrm{lim}}\nolimits}
\def\colim{\mathop{\mathrm{colim}}\nolimits}
\def\Spec{\mathop{\mathrm{Spec}}}
\def\Hom{\mathop{\mathrm{Hom}}\nolimits}
\def\Ext{\mathop{\mathrm{Ext}}\nolimits}
\def\SheafHom{\mathop{\mathcal{H}\!\mathit{om}}\nolimits}
\def\SheafExt{\mathop{\mathcal{E}\!\mathit{xt}}\nolimits}
\def\Sch{\mathit{Sch}}
\def\Mor{\mathop{\mathrm{Mor}}\nolimits}
\def\Ob{\mathop{\mathrm{Ob}}\nolimits}
\def\Sh{\mathop{\mathit{Sh}}\nolimits}
\def\NL{\mathop{N\!L}\nolimits}
\def\CH{\mathop{\mathrm{CH}}\nolimits}
\def\proetale{{pro\text{-}\acute{e}tale}}
\def\etale{{\acute{e}tale}}
\def\QCoh{\mathit{QCoh}}
\def\Ker{\mathop{\mathrm{Ker}}}
\def\Im{\mathop{\mathrm{Im}}}
\def\Coker{\mathop{\mathrm{Coker}}}
\def\Coim{\mathop{\mathrm{Coim}}}

% Boxtimes
%
\DeclareMathSymbol{\boxtimes}{\mathbin}{AMSa}{"02}

%
% Macros for moduli stacks/spaces
%
\def\QCohstack{\mathcal{QC}\!\mathit{oh}}
\def\Cohstack{\mathcal{C}\!\mathit{oh}}
\def\Spacesstack{\mathcal{S}\!\mathit{paces}}
\def\Quotfunctor{\mathrm{Quot}}
\def\Hilbfunctor{\mathrm{Hilb}}
\def\Curvesstack{\mathcal{C}\!\mathit{urves}}
\def\Polarizedstack{\mathcal{P}\!\mathit{olarized}}
\def\Complexesstack{\mathcal{C}\!\mathit{omplexes}}
% \Pic is the operator that assigns to X its picard group, usage \Pic(X)
% \Picardstack_{X/B} denotes the Picard stack of X over B
% \Picardfunctor_{X/B} denotes the Picard functor of X over B
\def\Pic{\mathop{\mathrm{Pic}}\nolimits}
\def\Picardstack{\mathcal{P}\!\mathit{ic}}
\def\Picardfunctor{\mathrm{Pic}}
\def\Deformationcategory{\mathcal{D}\!\mathit{ef}}
\newenvironment{reference}{}{}
\newenvironment{slogan}{}{}
\newenvironment{history}{}{}
\usepackage{cleveref}
\crefname{originalsection}{original section}{original sections}
\Crefname{originalsection}{Original section}{Original sections}
\makeatletter
\newlabel{situation-regular-model}{{0C61}{source}{}{sourcepointer-1}{}}
\newlabel{situation-regular-model@cref}{{[originalsection][0][]0C61}{[1][0][]source}{}{}{}}
\makeatother
\crefname{originallemma}{original lemma}{original lemmas}
\Crefname{originallemma}{Original lemma}{Original lemmas}
\makeatletter
\newlabel{lemma-intersection-pairing}{{0C64}{source}{}{sourcepointer-2}{}}
\newlabel{lemma-intersection-pairing@cref}{{[originallemma][0][]0C64}{[1][0][]source}{}{}{}}
\makeatother
\crefname{originallemma}{original lemma}{original lemmas}
\Crefname{originallemma}{Original lemma}{Original lemmas}
\makeatletter
\newlabel{lemma-properties-form}{{0C66}{source}{}{sourcepointer-3}{}}
\newlabel{lemma-properties-form@cref}{{[originallemma][0][]0C66}{[1][0][]source}{}{}{}}
\makeatother
\crefname{originallemma}{original lemma}{original lemmas}
\Crefname{originallemma}{Original lemma}{Original lemmas}
\makeatletter
\newlabel{lemma-regular-model-genus}{{0C69}{source}{}{sourcepointer-4}{}}
\newlabel{lemma-regular-model-genus@cref}{{[originallemma][0][]0C69}{[1][0][]source}{}{}{}}
\makeatother
\crefname{originallemma}{original lemma}{original lemmas}
\Crefname{originallemma}{Original lemma}{Original lemmas}
\makeatletter
\newlabel{divisors-lemma-sum-effective-Cartier-divisors-union}{{0C4R}{source}{}{sourcepointer-5}{}}
\newlabel{divisors-lemma-sum-effective-Cartier-divisors-union@cref}{{[originallemma][0][]0C4R}{[1][0][]source}{}{}{}}
\makeatother
\crefname{originallemma}{original lemma}{original lemmas}
\Crefname{originallemma}{Original lemma}{Original lemmas}
\makeatletter
\newlabel{morphisms-lemma-scheme-theoretic-union}{{0C4J}{source}{}{sourcepointer-6}{}}
\newlabel{morphisms-lemma-scheme-theoretic-union@cref}{{[originallemma][0][]0C4J}{[1][0][]source}{}{}{}}
\makeatother
\crefname{originallemma}{original lemma}{original lemmas}
\Crefname{originallemma}{Original lemma}{Original lemmas}
\makeatletter
\newlabel{coherent-lemma-cohomology-projective-space-over-ring}{{01XT}{source}{}{sourcepointer-7}{}}
\newlabel{coherent-lemma-cohomology-projective-space-over-ring@cref}{{[originallemma][0][]01XT}{[1][0][]source}{}{}{}}
\makeatother
\title{Two exceptional curves of the first kind on a regular model can meet only in the degenerate configuration of two rational curves and genus zero}
\author{The Stacks Project Authors}
\date{stacks-project:models.tex:local-snapshot-20261012; GNU FDL 1.2 or later}
\begin{document}
\maketitle
\noindent\textit{Editorial note.} The counted unit is the Stacks Project Lemma lemma-exceptional-curves-dont-meet at source lines 4927-4935 together with its proof at source lines 4937-4981, a straight adjacent pairing in the chapter file. The chapter is the upstream file models.tex; the pool directory carries the longer name stacks-models and is not used as the chapter name anywhere in this card.\par\medskip
\noindent\textit{Editorial note.} This card has no context span; the situation named by the statement is declared as a pointer with the kind section, because the assembler's whitelist of kinds has no entry for a situation.\par\medskip
\noindent\textit{Editorial note.} The retained proof names seven results: three lemmas of this chapter and four results of the Divisors, Morphisms and Cohomology of Schemes chapters. All are declared as non-load-bearing navigation pointers with their stable Stacks tags looked up in the shared tag index; none of their intervals is reproduced here.\par\medskip
\noindent\textit{Editorial note.} The retained proof span starts at line 4937, the line of the opening marker read from the source, and contains exactly one proof environment.\par\medskip
\noindent\textit{Editorial note.} No citation occurs in the retained spans, so no bibliography entry is supplied.\par\medskip
\noindent\textit{Allowed standard prerequisites (editorial).}\begin{itemize}
\item The regular-model situation, the degree pairing of the components and the genus of the curve. The situation fixes a discrete valuation ring, a smooth projective curve over its fraction field, a regular proper model and the components of the special fibre with their multiplicities; the degree pairing on combinations of components is symmetric and is computed from the degrees of the restrictions of the invertible sheaves attached to the components, and the genus of the special fibre can be read off from the numerical type and equals the genus of the curve. For a union of two effective Cartier divisors the structure sheaf fits into a short exact sequence comparing it with the sheaves of the parts, and the cohomology of the projective line over the residue field is known. The author's own lemmas of the chapter and the quoted results of the Divisors, Morphisms and Coherent chapters are the allowed inputs.. The counted lemma is the constraint that keeps two minus-one curves from meeting in a general configuration, and it is used when minimal models are compared through sequences of contractions.
\end{itemize}\medskip
\ifdefined\hypertarget\hypertarget{sourcepointer-1}{}\fi
\noindent\textit{Source reference.} Original section 0C61: focused-seeds/AGP/sources/stacks-snapshot-20261011/tags.txt, source lines 15281--15281. Pointer to the situation named by the retained statement, declared with the kind section because the assembler's whitelist of kinds has no entry for a situation. The number is the stable Stacks tag of the target, looked up in the shared tag index.\par\medskip
\ifdefined\hypertarget\hypertarget{sourcepointer-2}{}\fi
\noindent\textit{Source reference.} Original lemma 0C64: focused-seeds/AGP/sources/stacks-snapshot-20261011/tags.txt, source lines 15284--15284. Navigation pointer invoked by the retained proof for the symmetry of the degree pairing on combinations of components, which is what lets the intersection number of the two curves be computed; that lemma is counted by a separate card of this lane, so its proof is not reproduced here. The number is the stable Stacks tag of the target, looked up in the shared tag index.\par\medskip
\ifdefined\hypertarget\hypertarget{sourcepointer-3}{}\fi
\noindent\textit{Source reference.} Original lemma 0C66: focused-seeds/AGP/sources/stacks-snapshot-20261011/tags.txt, source lines 15286--15286. Navigation pointer invoked by the retained proof for the numerical properties of the intersection form; the cited statement is an allowed input of the same chapter and its proof is not reproduced here. The number is the stable Stacks tag of the target, looked up in the shared tag index.\par\medskip
\ifdefined\hypertarget\hypertarget{sourcepointer-4}{}\fi
\noindent\textit{Source reference.} Original lemma 0C69: focused-seeds/AGP/sources/stacks-snapshot-20261011/tags.txt, source lines 15289--15289. Navigation pointer invoked by the retained proof for the identification of the genus of the curve with the genus read off from the numerical type; the cited statement is an allowed input of the same section and its proof is not reproduced here. The number is the stable Stacks tag of the target, looked up in the shared tag index.\par\medskip
\ifdefined\hypertarget\hypertarget{sourcepointer-5}{}\fi
\noindent\textit{Source reference.} Original lemma 0C4R: focused-seeds/AGP/sources/stacks-snapshot-20261011/tags.txt, source lines 15236--15236. Allowed input of the Divisors chapter quoted by the retained proof for the structure of the sum of effective Cartier divisors as a union; the pointer is registered so the reference resolves and the cited statement is not reproduced inside the span. The number is the stable Stacks tag of the target, looked up in the shared tag index.\par\medskip
\ifdefined\hypertarget\hypertarget{sourcepointer-6}{}\fi
\noindent\textit{Source reference.} Original lemma 0C4J: focused-seeds/AGP/sources/stacks-snapshot-20261011/tags.txt, source lines 15229--15229. Allowed input of the Morphisms chapter quoted by the retained proof for the identification of the scheme-theoretic union of the two curves; the pointer is registered so the reference resolves and the cited statement is not reproduced inside the span. The number is the stable Stacks tag of the target, looked up in the shared tag index.\par\medskip
\ifdefined\hypertarget\hypertarget{sourcepointer-7}{}\fi
\noindent\textit{Source reference.} Original lemma 01XT: focused-seeds/AGP/sources/stacks-snapshot-20261011/tags.txt, source lines 2532--2532. Allowed input of the Cohomology of Schemes chapter quoted by the retained proof for the cohomology of the projective line over a ring; the pointer is registered so the reference resolves and the cited statement is not reproduced inside the span. The number is the stable Stacks tag of the target, looked up in the shared tag index.\par\medskip
\begin{lemma}
\label{lemma-exceptional-curves-dont-meet}
In Situation \ref{situation-regular-model} given a pair of indices $i, j$
such that $C_i$ and $C_j$ are exceptional curves of the first kind
and $C_i \cap C_j \not = \emptyset$, then
$n = 2$, $m_1 = m_2 = 1$, $C_1 \cong \mathbf{P}^1_k$,
$C_2 \cong \mathbf{P}^1_k$, $C_1$ and $C_2$ meet in a $k$-rational point,
and $C$ has genus $0$.
\end{lemma}

\begin{proof}
Choose isomorphisms $C_i = \mathbf{P}^1_{\kappa_i}$ and
$C_j = \mathbf{P}^1_{\kappa_j}$. The scheme $C_i \cap C_j$
is a nonempty effective Cartier divisor in both $C_i$ and $C_j$.
Hence
$$
(C_i \cdot C_j) = \deg(C_i \cap C_j) \geq \max([\kappa_i: k], [\kappa_j : k])
$$
The first equality was shown in the proof of
Lemma \ref{lemma-intersection-pairing}.
On the other hand, the self intersection $(C_i \cdot C_i)$ is equal
to the degree of $\mathcal{O}_X(C_i)$ on $C_i$ which is $-[\kappa_i : k]$
as $C_i$ is an exceptional curve of the first kind. Similarly for
$C_j$. By Lemma \ref{lemma-properties-form}
$$
0 \geq (C_i + C_j)^2 = -[\kappa_i : k] + 2(C_i \cdot C_j) - [\kappa_j : k]
$$
This implies that $[\kappa_i : k] = \deg(C_i \cap C_j) = [\kappa_j : k]$
and that we have $(C_i + C_j)^2 = 0$. Looking at the lemma again
we conclude that $n = 2$, $\{1, 2\} = \{i, j\}$, and $m_1 = m_2$.
Moreover, the scheme theoretic intersection $C_i \cap C_j$ consists of
a single point $p$ with residue field $\kappa$ and
$\kappa_i \to \kappa \leftarrow \kappa_j$ are isomorphisms.
Let $D = C_1 + C_2$ as effective Cartier divisor on $X$.
Observe that $D$ is the scheme theoretic union of $C_1$ and $C_2$
(Divisors, Lemma \ref{divisors-lemma-sum-effective-Cartier-divisors-union})
hence we have a short exact sequence
$$
0 \to \mathcal{O}_D \to \mathcal{O}_{C_1} \oplus \mathcal{O}_{C_2} \to
\mathcal{O}_p \to 0
$$
by Morphisms, Lemma \ref{morphisms-lemma-scheme-theoretic-union}.
Since we know the cohomology of $C_i \cong \mathbf{P}^1_\kappa$
(Cohomology of Schemes, Lemma
\ref{coherent-lemma-cohomology-projective-space-over-ring})
we conclude from the long exact cohomology sequence that
$H^0(D, \mathcal{O}_D) = \kappa$ and
$H^1(D, \mathcal{O}_D) = 0$. By Lemma \ref{lemma-regular-model-genus}
we conclude
$$
1 - g_C = d[\kappa : k](1 - 0)
$$
where $d = m_1 = m_2$. It follows that $g_C = 0$ and $d = m_1 = m_2 = 1$
and $\kappa = k$.
\end{proof}

\end{document}
